\documentclass[11pt]{article}

\usepackage[margin=1in]{geometry}
\usepackage{amsmath}
\usepackage{amssymb}
\usepackage{hyperref}

\title{Stage 3A Physics Note: Initial Data for the 5D Uniform Black String}
\author{GL-with-AI Project}
\date{\today}

\newcommand{\GRSP}{\texttt{GR\_SPACEDIM}}
\newcommand{\CHSP}{\texttt{CH\_SPACEDIM}}

\begin{document}
\maketitle

\section{Scope and Status}

This note is a polished physics review artifact for Stage 3A. It records the
target continuum initial data and design choices for a future implementation of
the five-dimensional uniform black string and a Gregory-Laflamme seed
perturbation.

This is design only. The initial data described here are not implemented in
\texttt{BlackStringToy}. The current \texttt{BlackStringToy} scaffold remains a
non-physical build/run scaffold derived from the public GRChombo
\texttt{Examples/BinaryBH} example. Nothing in this note claims that the
current code evolves a five-dimensional black string.

\section{Target Geometry}

The target unperturbed spacetime is the five-dimensional uniform black string,
formally
\begin{equation}
  \mathrm{Schwarzschild}_4 \times S^1_z .
\end{equation}

The dimensional bookkeeping for the intended reduced evolution is:
\begin{itemize}
  \item physical spacetime dimension: \(5\);
  \item physical spatial dimension: \(\GRSP = 4\);
  \item Chombo gridded dimension: \(\CHSP = 2\);
  \item gridded directions: a radial/cartoon coordinate \(x\) and compact
    string coordinate \(z\);
  \item hidden/reconstructed directions: the transverse \(S^2\), represented by
    SO(3)/modified-cartoon variables and source terms.
\end{itemize}

\section{GP-like Horizon-Penetrating Form}

The planned continuum starting point is a horizon-penetrating
Gullstrand-Painleve-like form of the uniform string. A schematic line element is
\begin{equation}
  ds^2
  =
  -dt^2
  +
  \left(dx + s\sqrt{\frac{r_0}{x}}\,dt\right)^2
  +
  x^2 d\Omega_2^2
  +
  dz^2 .
  \label{eq:gp-line-element}
\end{equation}
Here \(s=\pm 1\) labels the GP branch/sign convention. The slice is
horizon-penetrating: the horizon at \(x=r_0\) is not a coordinate singularity
on the spatial slice.

The sign \(s\), the time orientation, and GRChombo's convention for
extrinsic curvature \(K_{ij}\) must be matched before implementation. All
extrinsic-curvature formulas below should be read with that caveat.

\section{ADM Variables and GP Extrinsic Curvature}

The intrinsic spatial metric on the GP-like slice is flat in the radial
transverse sector times the compact string direction:
\begin{equation}
  dl^2 = dx^2 + x^2 d\Omega_2^2 + dz^2 .
\end{equation}

The natural GP lapse and shift associated with
Eq.~\eqref{eq:gp-line-element} are
\begin{equation}
  \alpha = 1,
  \qquad
  \beta^x = s\sqrt{\frac{r_0}{x}},
  \qquad
  \beta^z = 0 .
  \label{eq:natural-gp-shift}
\end{equation}

Using the convention
\begin{equation}
  K_{ij}
  =
  \frac{1}{2\alpha}
  \left(
    D_i\beta_j + D_j\beta_i - \partial_t \gamma_{ij}
  \right),
  \label{eq:kij-convention}
\end{equation}
and the time-independent spatial metric, the nonzero components for the
uniform string are, up to GRChombo's possible overall \(K_{ij}\) sign
convention,
\begin{align}
  K_{xx}
  &=
  -\frac{s}{2}\sqrt{\frac{r_0}{x^3}},
  \\
  K_{ww}
  &=
  s\sqrt{\frac{r_0}{x^3}}
  \qquad\text{for each hidden transverse direction},
  \\
  K_{zz}
  &= 0 .
\end{align}
Thus \(K_{zz}=0\), but the full four-spatial-dimensional trace is not zero.
The trace is
\begin{equation}
  K
  =
  K^x{}_{x}
  +
  K^z{}_{z}
  +
  K^{w_1}{}_{w_1}
  +
  K^{w_2}{}_{w_2}
  =
  \frac{3s}{2}\sqrt{\frac{r_0}{x^3}},
  \label{eq:k-trace}
\end{equation}
again up to the overall GRChombo \(K_{ij}\) sign convention. If GRChombo uses
the opposite convention to Eq.~\eqref{eq:kij-convention}, then all
\(K_{ij}\), \(K\), and conformal traceless components \(A_{ij}\) and
\(A_{ww}\) must be sign-flipped.

\section{Scalar and Sign Note for \(K\)}

The quantity \(K=\gamma^{ij}K_{ij}\) is a scalar on a fixed spatial slice, so
it is invariant under spatial coordinate changes on that slice. Its sign can
still change if one reverses the GP branch, reverses the time orientation, or
uses the opposite convention for the future-pointing normal. The two values
\(s=\pm 1\) are time-reversed ingoing/outgoing embeddings of the same flat
intrinsic spatial geometry. Therefore the sign of \(K\) is fixed only after
choosing both the GP branch and GRChombo's \(K_{ij}\) convention.

\section{Conformal CCZ4 Variable Map}

The target variables follow the conformal CCZ4 convention already recorded in
the implementation notes:
\begin{equation}
  \gamma_{ij} = \chi^{-1} h_{ij},
  \qquad
  \gamma_{ww} = \chi^{-1} h_{ww}.
  \label{eq:conformal-metric}
\end{equation}
The hidden physical extrinsic-curvature component is reconstructed as
\begin{equation}
  K_{ww}
  =
  \chi^{-1}
  \left(
    A_{ww} + \frac{h_{ww}K}{\GRSP}
  \right).
  \label{eq:kww-reconstruction}
\end{equation}
For \(\GRSP=4\),
\begin{equation}
  A_{ij}
  =
  \chi
  \left(
    K_{ij} - \frac{1}{4}\gamma_{ij}K
  \right),
  \qquad
  A_{ww}
  =
  \chi K_{ww} - \frac{1}{4}h_{ww}K .
  \label{eq:aij-map}
\end{equation}

At design level, the future initial-data class must provide:
\begin{itemize}
  \item \(\chi\), initially \(1\) before the GL perturbation;
  \item conformal metric components \(h_{xx}\), \(h_{zz}\), and \(h_{xz}\);
  \item hidden conformal metric component \(h_{ww}\);
  \item full trace \(K\), using all four spatial directions;
  \item traceless components \(A_{xx}\), \(A_{zz}\), \(A_{xz}\), and \(A_{ww}\);
  \item lapse, shift, Gamma variables, and gauge-driver variables consistent
    with the chosen startup convention.
\end{itemize}

For the unperturbed flat GP-like spatial slice, the expected conformal metric
starting values are
\begin{equation}
  h_{xx}=1,
  \qquad
  h_{zz}=1,
  \qquad
  h_{xz}=0,
  \qquad
  h_{ww}=1,
\end{equation}
before any GL perturbation is applied.

\section{Meaning of \(h_{ww}=1\)}

The variable \(h_{ww}\) is not the angular metric component
\(g_{\theta\theta}\). It is the hidden Cartesian-like conformal metric
component used by the SO(3)/modified-cartoon reduction. In flat transverse
space, the suppressed Cartesian directions have \(\gamma_{ww}=1\).

The angular factor \(x^2\) in \(x^2d\Omega_2^2\) is supplied by the
coordinate/cartoon reconstruction. The physical transverse sphere radius is
\begin{equation}
  R
  =
  x\sqrt{\gamma_{ww}}
  =
  x\sqrt{\chi^{-1}h_{ww}} .
  \label{eq:cartoon-radius}
\end{equation}
Setting \(h_{ww}=x^2\) would double-count the angular radius.

\section{Horizon Representation}

The future apparent-horizon representation is
\begin{equation}
  x = h(t,z).
\end{equation}
Here \(h(t,z)\) is the coordinate radius/location of the apparent horizon on
the reduced grid. For the unperturbed uniform string,
\begin{equation}
  h(0,z)=r_0 .
\end{equation}
During Gregory-Laflamme evolution, \(h(t,z)\) becomes \(z\)-dependent. The
physical transverse \(S^2\) radius on the horizon is
\begin{equation}
  R(t,z)
  =
  h(t,z)
  \sqrt{
    \gamma_{ww}\bigl(t,h(t,z),z\bigr)
  } .
  \label{eq:horizon-radius}
\end{equation}

\section{GL Seed Perturbation}

The planned Gregory-Laflamme seed from the reproduction roadmap is a
conformal-factor perturbation localized near the horizon:
\begin{equation}
  \chi
  =
  1
  +
  \epsilon
  \sin\left(\frac{2\pi n z}{L}\right)
  \exp\left[
    -
    \left(
      \frac{x}{r_0}
      -
      \frac{r_0}{x}
    \right)^2
  \right].
  \label{eq:gl-seed}
\end{equation}
Planned defaults are
\begin{equation}
  \epsilon=0.01,
  \qquad
  n=1,
  \qquad
  r_0=1,
  \qquad
  L/r_0=10 .
\end{equation}
This perturbation is planned but not implemented. It intentionally introduces
small constraint violations that must be monitored and damped before nonlinear
GL growth dominates.

\section{Natural GP Shift Versus Zero-Shift Gauge Startup}

The GP line element is used primarily to define a regular
horizon-penetrating initial hypersurface and to derive the geometric data
\(\gamma_{ij}\) and \(K_{ij}\). Natural GP coordinates have the nonzero shift
in Eq.~\eqref{eq:natural-gp-shift}.

The evolution gauge variables do not necessarily need to be initialized to the
natural GP shift. If following the reproduction roadmap, one can initialize
\begin{equation}
  \alpha(0)=1,
  \qquad
  \beta^i(0)=0,
\end{equation}
while keeping \(\gamma_{ij}\) and nonzero \(K_{ij}\) from the GP slice. In that
case the data are GP-slice geometric data, not a fully stationary
GP-coordinate evolution.

Do not set \(K_{ij}=0\) merely because the evolution shift is initialized to
zero. \(K_{ij}\) is part of the geometric initial data and must not be
recomputed from the chosen gauge-startup shift.

\section{Turduckening and Interior Handling}

The GP-like expressions include powers such as \(\sqrt{r_0/x}\), so a cutoff
or smoothing strategy is needed near \(x=0\). The planned cutoff scale is
\begin{equation}
  x_{\rm cut} = 0.1 r_0 .
\end{equation}
The quantities that require an interior prescription include \(\chi\), the GL
perturbation, the shift if used geometrically, \(K\), \(A_{ij}\), \(A_{ww}\),
\(h_{ww}\), and derived Gamma/constraint quantities. The exact smoothing
profile is unresolved. A hard clamp is simple but may create derivative
artifacts, so a smoother transition may be needed.

\section{Boundary and Periodicity}

The compact string direction should be periodic:
\begin{equation}
  z \sim z + L .
\end{equation}
The outer radial boundary should be far enough away that it is not causally
connected to the string during the target evolution window. The reproduction
roadmap uses \(x_{\rm outer}=256 r_0\). The long-term outer boundary will need
a reviewed outgoing or Sommerfeld-like treatment.

The cartoon axis and radial-origin regularity are separate issues. The
inherited BinaryBH reflective/Sommerfeld choices, binary-puncture tracking,
puncture-centered tagging, and Weyl4 extraction settings are not validated
long-term assumptions for the reduced 5D black-string system.

\section{Unresolved Decisions Before Coding}

\begin{itemize}
  \item Choose the GP branch \(s\) and match it to GRChombo's \(K_{ij}\) sign
    convention.
  \item Confirm whether to follow the roadmap startup
    \(\alpha(0)=1\), \(\beta^i(0)=0\) while using GP-slice
    \(\gamma_{ij}\) and nonzero \(K_{ij}\).
  \item Decide and verify the exact \texttt{hww}/\texttt{Aww} enum layout,
    including the AH convention \(\texttt{c\_hww}=\texttt{c\_K}-1\).
  \item Establish a reliable \(\GRSP=4\), \(\CHSP=2\) build path.
  \item Define unperturbed uniform-string Hamiltonian and momentum constraint
    validation tests.
  \item Compare horizon radius and area against \(r_0\) and
    \(A_0 = 4\pi r_0^2 L\).
  \item Choose and validate the turduckening smoothing profile.
\end{itemize}

\end{document}
